\begin{definition}\label{sga3:XXI:2.1.3}
A root \(\alpha\in R\) is called \emph{indivisible} if
\(\alpha/2\notin R\). The root datum \(\cR\) is called
\emph{reduced} if every root is indivisible.
\end{definition}

The root datum \(\cR\) is reduced if and only if \(\cR^\ast\) is
reduced. If \(\alpha\) is indivisible and \(w\in W\), then
\(w(\alpha)\) is indivisible.

\begin{definition}\label{sga3:XXI:2.1.4}
Let \(\alpha\in R\). If \(\alpha\) is indivisible, set
\(\operatorname{ind}(\alpha)=\alpha\); otherwise set
\(\operatorname{ind}(\alpha)=\alpha/2\).
\end{definition}

\begin{corollary}\label{sga3:XXI:2.1.5}
If \(\alpha\in R\) is indivisible and \(k\alpha\in R\), where
\(k\in\mathbb Q\), then \(k\in\mathbb Z\).
\end{corollary}

\begin{proposition}\label{sga3:XXI:2.1.6}
Let \(\cR\) be a root datum. Then
\[
 \operatorname{ind}(\cR)
 =\bigl(M,M^\ast,\operatorname{ind}(R),
         \operatorname{ind}(R)^\ast\bigr)
\]
is a reduced root datum, and
\[
 W\bigl(\operatorname{ind}(\cR)\bigr)=W(\cR).
\]
\end{proposition}

\begin{proof}
The quadruple \(\operatorname{ind}(\cR)\) is a root datum by
Corollary~\ref{sga3:XXI:1.2.11}, since the Weyl group permutes the
indivisible roots. The second assertion follows from
Proposition~\ref{sga3:XXI:2.1.1}.
\end{proof}

\begin{remark}\label{sga3:XXI:2.1.7}
If \(\cR\) is not reduced, then
\[
 \operatorname{ind}(R)^\ast\neq\operatorname{ind}(R^\ast),
 \qquad
 \operatorname{ind}(\cR^\ast)\neq\operatorname{ind}(\cR)^\ast.
\]
\end{remark}

\subsection{Orthogonal roots}

\begin{lemma}\label{sga3:XXI:2.2.1}
Let \(\alpha\) and \(\beta\) be two roots. Then
\begin{equation}
 \ell(\alpha)(\alpha^\ast,\beta)
 =\ell(\beta)(\beta^\ast,\alpha).
 \tag{11}\label{sga3:XXI:eq:11}
\end{equation}
\end{lemma}

\begin{proof}
This follows immediately from Lemma~\ref{sga3:XXI:1.2.1},
formulas \eqref{sga3:XXI:eq:9} and \eqref{sga3:XXI:eq:10}.
\end{proof}

\begin{corollary}\label{sga3:XXI:2.2.2}
Let \(\alpha,\beta\in R\). The following conditions are equivalent:
\begin{enumerate}
\item[\SGAInlineRefTarget{sga3:XXI:2.2.2:item:i}(i)] \((\alpha^\ast,\beta)=0\);
\item[(i bis)] \((\beta^\ast,\alpha)=0\);
\item[(ii)] \((p(\alpha),\beta)=0\);
\item[(iii)] \(s_\alpha(\beta)=\beta\);
\item[(iii bis)] \(s_\beta(\alpha)=\alpha\);
\item[(iv)] \(s_\alpha(\beta^\ast)=\beta^\ast\);
\item[(iv bis)] \(s_\beta(\alpha^\ast)=\alpha^\ast\);
\item[\SGAInlineRefTarget{sga3:XXI:2.2.2:item:v}(v)] \(s_\alpha\neq s_\beta\), and \(s_\alpha\) and \(s_\beta\)
  commute.
\end{enumerate}
\end{corollary}

\begin{proof}
All the equivalences are immediate except those involving \SGARefLink{sga3:XXI:2.2.2:item:v}{(v)}. We
first show that \SGARefLink{sga3:XXI:2.2.2:item:i}{(i)}, together with its equivalent condition (i bis),
implies \SGARefLink{sga3:XXI:2.2.2:item:v}{(v)}. One has
\[
 s_\alpha s_\beta(x)
 =x-(\beta^\ast,x)\beta-(\alpha^\ast,x)\alpha
   +(\beta^\ast,x)(\alpha^\ast,\beta)\alpha.
\]
If \((\alpha^\ast,\beta)=0\), then
\((\beta^\ast,\alpha)=0\), and consequently
\[
 s_\alpha s_\beta(x)=s_\beta s_\alpha(x).
\]

Conversely, suppose \SGARefLink{sga3:XXI:2.2.2:item:v}{(v)}. Then
\[
 s_\alpha=s_\beta s_\alpha s_\beta
          =s_{s_\beta(\alpha)}
 \qquad\text{by Corollary~\ref{sga3:XXI:1.2.10}.}
\]
By Proposition~\ref{sga3:XXI:2.1.1}, there is a
\(k\in\mathbb Q^\ast\) such that
\[
 \alpha=k\,s_\beta(\alpha)
       =k\alpha-k(\beta^\ast,\alpha)\beta.
\]
Since \(s_\alpha\neq s_\beta\), the roots \(\alpha\) and \(\beta\)
are not proportional, again by Proposition~\ref{sga3:XXI:2.1.1}.
It follows that \((\beta^\ast,\alpha)=0\).
\end{proof}

\begin{definition}\label{sga3:XXI:2.2.3}
Two roots satisfying the equivalent conditions of
Corollary~\ref{sga3:XXI:2.2.2} are called \emph{orthogonal}.
\end{definition}
